\chapter{الآلة كاملة}
\chaptermark{الآلة كاملة}
\label{sec:synthesis}

\section{ما الذي جمعناه}

أخذنا أنواع العصبونات في الجدول~\ref{tab:types} واحدًا تلو الآخر، وسألنا: ماذا يضيف كلٌّ منها إلى آلة حاسبة؟ وكانت القواعد في كل مرة واحدة: لا شيء إلا ما يوجد في كائن حي، وكل شيء يجري في زمن متصل. والآن نجمع الأجزاء في آلة واحدة ونرى كيف تعمل معًا.

الصورة التي رسمناها في الفصل~\ref{sec:neurontypes} صمدت وازدادت دقة. شبكة معنونة هائلة من عصبونات \nt{GLUT} تحمل البيانات. وعصبونات \nt{GABA} تتحكم في تدفق هذه البيانات. وفوق الشبكة أربع قنوات بثّ، لكلٍّ منها مقبضها: \nt{DOPA} يقول أيّ تغييرات الأوزان تُحفظ، و\nt{SERO} يمسك المستوى العام ويحدد، بحسب فرضية، أفق التخطيط، و\nt{NORA} يختار بين البحث والحسم، و\nt{ACET} بين الكتابة والقراءة (الشكل~\ref{fig:machine}).

\begin{figure}[t]
\centering
\begin{tikzpicture}[>=Stealth,
  blk/.style={draw,very thick,rounded corners=3pt,align=center,font=\small},
  reg/.style={draw,very thick,double,double distance=1.2pt,circle,minimum size=1.25cm,inner sep=0pt,font=\scriptsize,fill=orange!15},
  bc/.style={->,thick,densely dashed,orange!85!black}]
% sensors, bus, actions
\node[blk,fill=green!8,text width=1.7cm] (S) at (0.9,0) {المستشعرات\\\scriptsize تشغيل/إطفاء};
\draw[very thick,rounded corners=4pt,fill=blue!5] (2.6,-1.35) rectangle (9.0,1.35);
\node[font=\small] at (5.8,0.95) {ناقل \nt{GLUT}: البيانات};
\node[font=\scriptsize,align=center] at (4.3,0.25) {التزامنات، التأخيرات،\\المنطق (الفصل~\ref{sec:glutcomputer})};
\node[font=\scriptsize,align=center] at (7.4,0.25) {الذاكرة، الشفرة\\(الفصلان~\ref{sec:code} و\ref{sec:acet})};
\draw[thick,red!70!black,fill=red!8,rounded corners=2pt] (2.8,-1.15) rectangle (8.8,-0.55);
\node[font=\scriptsize,red!60!black] at (5.8,-0.85) {\nt{GABA} (الفصل~\ref{sec:gaba}): منع، اختيار، قسمة، إيقاع};
\node[blk,fill=blue!5,text width=2.0cm] (A) at (10.9,0) {اختيار\\الفعل};
\draw[->,very thick] (S) -- (2.6,0);
\draw[->,very thick] (9.0,0) -- (A);
\draw[->,very thick] (A) -- (12.6,0) node[right,font=\small] {العضلات};
\pic[scale=1.1] at (13.2,-0.5) {spring};
\pic[scale=1.15] at (0.4,0.85) {eyepic};
\pic[scale=1.05] at (1.35,0.85) {speaker};
% registers
\node[reg] (D) at (3.4,3.2) {\nt{DOPA}};
\node[reg] (C) at (5.6,3.2) {\nt{SERO}};
\node[reg] (N) at (7.8,3.2) {\nt{NORA}};
\node[reg] (H) at (10.0,3.2) {\nt{ACET}};
\node[font=\scriptsize,align=center,above=1pt] at (D.north) {الخطأ $\delta$};
\node[font=\scriptsize,align=center,above=1pt] at (C.north) {المستوى، $\tau_\gamma$};
\node[font=\scriptsize,align=center,above=1pt] at (N.north) {الكسب $g$};
\node[font=\scriptsize,align=center,above=1pt] at (H.north) {الكتابة $a$};
\draw[bc] (D.south) -- (3.4,1.35);
\draw[bc] (C.south) -- (5.6,1.35);
\draw[bc] (N.south) -- (7.8,1.35);
\draw[bc] (H.south) -- (10.0,2.1) -- (8.8,1.35);
\draw[bc] (H.south east) to[bend left=20] (A.north east);
\draw[->,thick] (2.85,1.35) -- (2.85,2.75) -- (D.west) node[pos=0.25,left,font=\scriptsize] {$V$};
\draw[->,thick,green!45!black] (0.4,3.2) node[left,black,font=\small] {المكافأة $r$} -- (2.2,3.2) -- (D.west);
\pic[scale=0.9] at (-0.45,3.7) {juice};
\draw[->,thick,densely dotted,gray!50!black] ([yshift=0.45cm]D.north) to[bend left=18] node[above,font=\scriptsize,black] {المفاجأة} ([yshift=0.45cm]N.north);
\end{tikzpicture}
\caption{الآلة بكاملها. ينقل ناقل \nt{GLUT} البيانات من المستشعرات إلى اختيار الفعل، ويتحكم \nt{GABA} في التدفق داخله. الدوائر المزدوجة قنوات البثّ؛ والأسهم المتقطعة بثُّ إشاراتها إلى الشبكة كلها، ولكل قناة مقبضها. يتلقى \nt{DOPA} المكافأة $r$ والتوقع $V$ من الناقد داخل الناقل (الفصل~\ref{sec:dofa})؛ والسهم المنقّط فرضيةُ أن \nt{NORA} يعرف بالمفاجأة من أخطاء كبيرة على غير العادة (الفصل~\ref{sec:nora}).}
\label{fig:machine}
\end{figure}

\section{صيغة واحدة لكل المقابض}

يمكن جمع المقابض في صيغة تخطيطية واحدة. ليست هذه نموذجًا جديدًا، بل خلاصة نماذج الفصول من~\ref{sec:glutcomputer} إلى~\ref{sec:acet}: كل رمز مأخوذ من فصله.
\begin{empheq}[box=\keybox]{equation}
\begin{aligned}
\tau\,\frac{\dd r_i}{\dd t} \;&=\; -r_i + F\Bigl(g\,\Bigl[\tfrac{1+a}{2}\,x_i + (1-a)\sum_j W_{ij}\,r_j - \sum_k M_{ik}\,q_k - \theta\Bigr]\Bigr),\\
\frac{\dd W_{ij}}{\dd t} \;&=\; \eta\,a\,\delta(t)\,e_{ij}(t),
\qquad \tau_e\,\frac{\dd e_{ij}}{\dd t} = -e_{ij} + r_i\,r_j,
\qquad \delta = r + \frac{\dd V}{\dd t} - \frac{V}{\tau_\gamma} .
\end{aligned}
\label{eq:machine}
\end{empheq}
هنا $r_i$ ترددات عصبونات \nt{GLUT}، و$x_i$ الدخل من المستشعرات، و$W_{ij}\ge0$ أوزانها المتبادلة (الفصل~\ref{sec:glutcomputer})؛ و$q_k$ ترددات عصبونات \nt{GABA}، و$M_{ik}\ge0$ شدة تثبيطها (الفصل~\ref{sec:gaba})؛ و$e_{ij}$ أثر الأهلية، و$\delta$ خطأ \nt{DOPA} (الفصل~\ref{sec:dofa})؛ و$\tau_\gamma$ الأفق الذي يحدده \nt{SERO} بحسب الفرضية؛ و$g$ كسب \nt{NORA} (الفصل~\ref{sec:nora})؛ و$a$ مستوى \nt{ACET} (الفصل~\ref{sec:acet}).

انظر إلى ما ليس موجودًا في~\eqref{eq:machine}. لا نبضة ساعة: كل شيء مكتوب بمعادلات تفاضلية، وكل عصبون يتغير من تلقاء نفسه. لا ذاكرة منفصلة: تتذكر الأوزان $W_{ij}$ نفسها التي يجري عليها الحساب، ويتذكر النشاط $r_i$ نفسه. لا تصحيحات فردية للغالبية العظمى من المشابك: تعلّمها حاصل ضرب أثر محلي في عدد واحد مشترك؛ ولا يوجد سلك خطأ مستقل إلى كل مخرج إلا في الدارة التي تتعلم الحركات (الفصل~\ref{sec:motorlearning}). وإشارات التحكم أربعة أنواع فقط، $\delta$ و$\tau_\gamma$ و$g$ و$a$، لثمانين مليار عصبون. وكلٌّ منها في الحقيقة ليس عددًا واحدًا، بل حزمة صغيرة من الخطوط المتوازية (القضية~\ref{prop:bundle})، لكن خطوط الحزمة تُعدّ بالآحاد.

\begin{keyblock}
\begin{proposition}[البنية]
\label{prop:architecture}
يمكن وصف آلة الدماغ، بقدر ما نفهمها اليوم، بثلاث طبقات. تسير البيانات على ناقل العناوين \nt{GLUT}. وتوجّه عصبونات \nt{GABA} المعنونة تدفقَ البيانات وتحدّه وتطبّعه. ويحدد نمطَ العمل عددٌ قليل من قنوات البث، كلٌّ منها حزمة صغيرة من الخطوط المتوازية المشتركة بين مناطق واسعة من الشبكة. الطبقتان الأوليان راسختان؛ أما توزيع الأدوار بين قنوات البث ففرضية في معظمه.
\end{proposition}
\end{keyblock}

\section{كل الأنواع في جدول واحد}

يجمع الجدول~\ref{tab:synthesis} كل أنواع العصبونات في الكتاب. الأنواع الأربعة التي لم يُخصَّص لها فصل نال كلٌّ منها سطرًا فيه؛ ووجد اثنان منها دورًا في فصول مخارج الآلة: \nt{GLYC} في حلقات الحبل الشوكي (الفصل~\ref{sec:muscles})، و\nt{ADRE} في المخرج الكيميائي (الفصل~\ref{sec:chemoutput}). أما دور الاثنين الآخرين في الحساب فإما بسيط وإما مجهول. والمواد التي لا تحدد نوع العصبون مجموعة في الملحق~\ref{sec:othermediators}.

\begin{table}[t]
\centering
\footnotesize
\setlength{\tabcolsep}{3pt}
\renewcommand{\arraystretch}{1.15}
\begin{tabular}{>{\raggedright}p{1.3cm}>{\raggedright}p{1.0cm}>{\raggedright}p{3.9cm}>{\raggedright}p{1.5cm}>{\raggedright}p{1.8cm}>{\raggedright\arraybackslash}p{3.8cm}}
\toprule
النوع & الفصل & ماذا يحسب & الزمن & الناقل & ما هو معروف \\
\midrule
\nt{GLUT} & \ref{sec:glutcomputer}, \ref{sec:code}, \ref{sec:sensors}, \ref{sec:motorlearning}, \ref{sec:dendrites} & البيانات: التزامنات، التأخيرات، المنطق، الذاكرة، التتابعات & \foreignlanguage{english}{ms} & معنون & راسخ \\
\nt{GABA} & \ref{sec:gaba}, \ref{sec:code}, \ref{sec:pain}, \ref{sec:motorlearning} & التحكم في التدفق: المنع، الاختيار، القسمة، الاستقرار، الإيقاع؛ بوابة الألم & \foreignlanguage{english}{ms} & معنون & راسخ؛ قسمة التردد بمعونة الضجيج الخلفي \\
\nt{SERO} & \ref{sec:serocomputer}, \ref{sec:pain}, \ref{sec:sleep} & منظّم المستوى، التنعيم؛ الأفق $\tau_\gamma$؛ مقبض صوت الألم؛ أنماط النوم & ثوانٍ & حجمي & التثبيط الذاتي مُقاس؛ الأفق فرضية \\
\nt{DOPA} & \ref{sec:dofa}, \ref{sec:dofaalt} & خطأ التنبؤ $\delta$؛ المستوى البطيء ثمنُ الوقت & $0.1\,\text{s}$ ودقائق & بثّ & $\delta$ مُقاس؛ التوسعات في الفصل~\ref{sec:dofaalt} \\
\nt{NORA} & \ref{sec:nora}, \ref{sec:pain}, \ref{sec:chemoutput}, \ref{sec:sleep} & الكسب $g$: البحث أم الحسم؛ المفاجأة؛ ”دوّاسة الوقود“ للأعضاء & ثوانٍ & بثّ & ارتفاع نسبة الإشارة إلى الضجيج مُقاس؛ الدور في البحث فرضية \\
\nt{ACET} & \ref{sec:acet}, \ref{sec:muscles}, \ref{sec:chemoutput}, \ref{sec:sleep} & الكتابة أو القراءة؛ سرعة التعلم؛ المخرج إلى العضلة؛ ”المكابح“ للأعضاء & ثوانٍ & كلاهما & ثلاثة تأثيرات مُقاسة؛ تبديل الأنماط فرضية \\
\nt{GLYC} & \ref{sec:muscles}, \ref{sec:pain} & وزن سالب في الحبل الشوكي وجذع الدماغ: منع العضلة المضادة & \foreignlanguage{english}{ms} & معنون & راسخ \\
\nt{HIST} & --- & إشارة عامة ”ابقَ يقظًا“ & بطيء & بثّ & دوره في الحساب لم يُدرس \\
\nt{ADRE} & \ref{sec:chemoutput} & ”دوّاسة الوقود“ للجسم كله عبر الدم؛ في الدماغ مجهول & بطيء & بثّ & خارج الدماغ راسخ؛ دوره في الدماغ غير واضح \\
\nt{ASPA} & --- & وزن موجب ضعيف & \foreignlanguage{english}{ms} & معنون & دوره موضع خلاف \\
\bottomrule
\end{tabular}
\caption{كل أنواع العصبونات في الكتاب: ماذا يحسب كلٌّ منها، وإلى أي حد نعرف ذلك.}
\label{tab:synthesis}
\end{table}

\section{كيف تعمل المقابض معًا}

حتى الآن كان كل فصل يدير مقبضًا واحدًا. فلنتتبع الآن حدثًا واحدًا عبر الآلة كلها (الشكل~\ref{fig:timeline}). تعلّم الحيوان منذ زمن أن الفعل $A$ يجلب العصير. وفي لحظة ما يتغير العالم: يكفّ $A$ عن جلب العصير، ويجلبه بدلًا منه الفعل $B$ الذي لا يكاد الحيوان يجربه.

\emph{الخطأ.} لم يأتِ العصير بعد $A$، فتجيب عصبونات \nt{DOPA} بهبوط، $\delta<0$ بحسب صيغة الخطأ~\eqref{eq:ctd}. والمشابك التي اختارت $A$ لتوّها تحمل أثرًا، فتضعف.

\emph{المفاجأة.} هبوط واحد مصادفة عادية. لكن الهبوطات تتكرر، وتصبح الأخطاء أكبر من المعتاد. وبحسب فرضية الفصل~\ref{sec:nora}، هذه بالذات إشارة \nt{NORA}: لقد تغيّر العالم. ينخفض الكسب $g$، ويصبح الاختيار ليّنًا، ويقود الضجيج إلى تجربة $B$ أكثر.

\emph{الكتابة.} بحسب فرضية الفصل~\ref{sec:acet}، يرتفع \nt{ACET} بعد التغيير وينقل الشبكة إلى نمط الكتابة: تضعف الوصلات الداخلية التي تحفظ العادة القديمة، ويقوى الدخل من المستشعرات، وتتغير الأوزان بسهولة.

\emph{عادة جديدة.} تجربة $B$ تجلب عصيرًا لم يتوقعه أحد: دفقة، $\delta>0$، فتقوى المشابك التي اختارت $B$. بضع محاولات كهذه، ويلحق التوقع $V$ بالعالم الجديد، وتعود الأخطاء صغيرة.

\emph{العودة.} مرّت المفاجأة، فيعيد \nt{NORA} الكسب العالي، ويعود الاختيار حادًّا؛ وينخفض \nt{ACET}، فتستعمل الشبكة في نمط القراءة ما تعلمته. وطوال هذا الوقت يحدد \nt{SERO}، بحسب الفرضية، الأفق $\tau_\gamma$: إلى أي مدى يستحق عصيرٌ بعيد أن يُنتظر.

\begin{figure}[t]
\centering
\begin{tikzpicture}[>=Stealth,x=1.1cm,y=0.95cm]
\foreach \y/\lab in {4/{$\delta$ (\nt{DOPA})},3/{$g$ (\nt{NORA})},2/{$a$ (\nt{ACET})},1/{اختيار $B$}}{
  \draw[gray!60!black] (0,\y-0.35) -- (10,\y-0.35);
  \node[left,font=\scriptsize] at (0,\y) {\lab};}
\draw[densely dashed,gray!60!black] (2,0.4) -- (2,4.55) node[above,font=\scriptsize,black] {تغيّر العالم};
\draw[densely dashed,gray!60!black] (5,0.4) -- (5,4.55) node[above,font=\scriptsize,black] {أول عصير مقابل $B$};
\pic[scale=0.85] at (2,5.25) {nojuice};
\pic[scale=0.85] at (5,5.25) {juice};
% delta: dips after the change, burst at the first juice for B, then back to zero
\draw[thick,red!70!black] (0,3.65) -- (2.3,3.65) -- (2.3,3.3) -- (2.5,3.3) -- (2.5,3.65) -- (3.2,3.65) -- (3.2,3.3) -- (3.4,3.3) -- (3.4,3.65) -- (4.1,3.65) -- (4.1,3.35) -- (4.3,3.35) -- (4.3,3.65) -- (5.0,3.65) -- (5.0,4.35) -- (5.2,4.35) -- (5.2,3.65) -- (5.8,3.65) -- (5.8,4.1) -- (6.0,4.1) -- (6.0,3.65) -- (6.6,3.65) -- (6.6,3.85) -- (6.8,3.85) -- (6.8,3.65) -- (10,3.65);
% gain: high, drops after surprise, recovers
\draw[thick,blue!60!black] (0,3.2) -- (3.0,3.2) -- (3.4,2.75) -- (5.8,2.75) -- (7.2,3.2) -- (10,3.2);
% ACh: low, rises after the change, falls after learning
\draw[thick,green!45!black] (0,1.75) -- (2.8,1.75) -- (3.3,2.25) -- (6.8,2.25) -- (7.8,1.75) -- (10,1.75);
% choice of B
\draw[thick,black] (0,0.7) -- (3.0,0.7) plot[domain=3:10,samples=40] (\x,{0.7+0.6/(1+exp(-(\x-5.3)*1.8))});
\draw[->,gray!70!black] (0,0.2) -- (10.2,0.2) node[below left,font=\scriptsize,black] {الزمن};
\end{tikzpicture}
\caption{حدث واحد متتبَّع عبر الآلة كلها. هذا مخطط نوعي مبني على نماذج الفصول من~\ref{sec:dofa} إلى~\ref{sec:acet} وفرضياتها، لا نتيجة حساب. بعد التغيير يجيب \nt{DOPA} بهبوطات؛ والهبوطات المتكررة، بحسب الفرضية، ترفع إشارة المفاجأة، فيخفض \nt{NORA} الكسب، ويشغّل \nt{ACET} الكتابة؛ ويعطي أول عصير مقابل $B$ دفقة، فينتقل الاختيار إلى $B$، وتعود المقابض إلى مواضعها.}
\label{fig:timeline}
\end{figure}

لاحظ أن أيّ مقبض لا يعرف ما تفعله المقابض الأخرى. كلٌّ منها يستجيب لعلاماته الخاصة، أي الخطأ وحجمه غير المعتاد وعدم اليقين، ويتكوّن ترتيب الأفعال من تلقاء نفسه، كما في دارة غير متزامنة ينطلق فيها كل مكوّن حين تجهز مداخله. أما هذا العمل المنسَّق بكامله فلم يُختبر بعد على حد علمنا: المقابض مقيسة كلٌّ على حدة، وعملها المشترك فرضية.

\section{بمَ تختلف هذه الآلة عن الحاسوب}

\begin{table}[t]
\centering
\footnotesize
\setlength{\tabcolsep}{4pt}
\renewcommand{\arraystretch}{1.15}
\begin{tabular}{>{\raggedright}p{2.2cm}>{\raggedright}p{4.2cm}>{\raggedright\arraybackslash}p{6.6cm}}
\toprule
& الحاسوب العادي & الدماغ \\
\midrule
الزمن & ساعة مشتركة، نحو جيجاهرتز & لا ساعة؛ كل عصبون يتغير بنفسه، والنوافذ تحددها الأحداث والإيقاعات المحلية (الفصول~\ref{sec:glutcomputer}، \ref{sec:gaba}، \ref{sec:code}) \\
العناصر & بوابات ثنائية موثوقة & عناصر عتبية تماثلية؛ يفقد المشبك معظم النبضات (الفصل~\ref{sec:code}) \\
علامة الوصلة & أيّ علامة & يحددها العصبون المرسل (الفصل~\ref{sec:neurontypes}) \\
الذاكرة & منفصلة عن المعالج & في الوصلات نفسها التي يجري فيها الحساب، وفي النشاط الذي يديم نفسه (الفصلان~\ref{sec:glutcomputer} و\ref{sec:acet}) \\
التعلّم & الانتشار العكسي للخطأ & قواعد محلية، وعدد بثّ واحد $\delta$ (الفصل~\ref{sec:dofa})، وأسلاك خطأ إلى مخارج دارة الحركة (الفصل~\ref{sec:motorlearning}) \\
التحكم & سجلات وناقل أوامر & أربع قنوات بثّ، كلٌّ منها حزمة من بضعة خطوط (الفصول~\ref{sec:serocomputer}، \babelsublr{\ref{sec:dofa}--\ref{sec:acet}}) \\
الطاقة & عشرات إلى مئات الواطات للمعالج الواحد & نحو $20$ واطًا للدماغ كله، بمعدل نبضة في الثانية تقريبًا لكل عصبون (الفصل~\ref{sec:code}) \\
\bottomrule
\end{tabular}
\caption{الدماغ والحاسوب العادي.}
\label{tab:computer}
\end{table}

يلخّص الجدول~\ref{tab:computer} الفروق الرئيسية. وليس أيٌّ منها عيبًا لم تجد الطبيعة وقتًا لإصلاحه، بل جزءًا من حل واحد. من دون ساعة مشتركة تلزم مستشعرات ”التشغيل والإطفاء“ وكواشف التزامن. والعناصر غير الموثوقة تحتاج إلى فائض من العصبونات الكثيرة. والذاكرة المدمجة في الحساب تحتاج إلى \nt{ACET} كي لا تفسد الكتابةُ القراءةَ. والتعلم بلا أسلاك راجعة يحتاج إلى \nt{DOPA} وإلى الضجيج. وحدّ الطاقة البالغ نبضة واحدة في الثانية يجبر اللغة كلها على التقتير وعلى إنفاق النبضات على الجديد وحده.

\section{ما لا نعرفه}

أصدق ما تُختم به هذه الخلاصة قائمةٌ بما لا نعرفه. وكل بند فيها مسألة لمهندس.

\emph{الشفرة.} نعرف سعة قناة النبضات، ونعرف أن المتلقي يختار أيّ جزء من الرسالة يقرأ (الفصل~\ref{sec:code}). لكن إلى أي حد تستعمل القشرة التوقيت الدقيق للنبضات، وهل للعصبونات ”كلمات“، فأمر مجهول.

\emph{التعلم عبر طبقات كثيرة.} إشارة البث تعلّم ببطء، وتزداد بطئًا مع كل عصبون يؤثر في النتيجة (القضية~\ref{prop:broadcastcost}). فكيف تضبط القشرة إذن مليارات الأوزان في دارات متعددة الطبقات؟ لا نعرف؛ وثمة نماذج تنقل فيها رشقاتُ النبضات الخطأَ بين الطبقات~\cite{Payeur2021}. وربما يكمن جزء من الجواب في الحسابات داخل فروع العصبون نفسه، حيث يتصرف عصبون واحد كشبكة صغيرة متعددة الطبقات: لتكرار استجابة نموذج عصبون قشري واحد تحتاج الشبكة الاصطناعية إلى خمس إلى ثماني طبقات~\cite{Beniaguev2021}، وفروع عصبونات الإنسان قادرة حتى على حساب ”أو الحصرية“~\cite{Gidon2020}. ومرشح آخر هو التعلم الذاتي: إذا تعلمت القشرة أن تتنبأ بمداخلها، نشأ الخطأ في موضعه، في كل طبقة، ولم تعد هناك حاجة إلى إشارة مشتركة (الفصل~\ref{sec:dofa})~\cite{Keller2018}.

\emph{ما الذي تنقله قنوات البث حقًّا.} لـ\nt{DOPA} صورة معيارية وستة توسعات لها (الفصل~\ref{sec:dofaalt})؛ ولـ\nt{NORA} و\nt{ACET} فرضيات معقولة (الفصلان~\ref{sec:nora} و\ref{sec:acet})؛ أما \nt{SERO} فمعظم ما لدينا عنه تخمينات.

\emph{التنسيق في الزمن.} رأينا كيف نحسب دالة، وكيف نقيس الزمن من حدث، وكيف نقطّع التدفق إلى نوافذ بإيقاع محلي. لكن كيف تنسّق شبكة هائلة بلا ساعة مشتركة عمل مناطق متباعدة، فأمر لا نفهمه إلا جزئيًا.

\emph{الطاقة.} عشرون واطًا لكل شيء. نفهم لماذا يجب أن تكون الشفرة متناثرة (القضية~\ref{prop:economy})، لكن لا أحد يعرف حتى الآن كيف يبني آلة تفعل الشيء نفسه بالقدرة نفسها~\cite{MehonicKenyon2022}.

الدماغ آلة حاسبة لم يجمعها مهندس، لكنها مبنية وفق قوانين يعرفها أي مهندس: دارات \foreignlanguage{english}{RC}، والعتبات، والتغذية الراجعة، والمرشحات، والنواقل، وسجلات البث. لم نقرأ مخططها إلا جزئيًا. وسيكمل قراءته من يجيد قراءة المخططات.

\section{ما التالي في الكتاب}

جمع هذا الفصل النواة الحسابية للآلة. أما بقية الفصول فتخرج عن حدودها. الفصلان~\ref{sec:sensors} و\ref{sec:recognition} عن المداخل: كيف تحوّل المستشعرات الضوء والصوت والجزيئات إلى نبضات، وكيف تتعرّف الآلة فيها على الأشياء. والفصول من~\ref{sec:muscles} إلى~\ref{sec:chemoutput} عن المخارج وما يخدمها: العضلات، والألم بوصفه إشارة إنذار، وتعلّم الحركات عبر أسلاك خطأ مستقلة، والمخرج الكيميائي البطيء إلى الجسم. والفصل~\ref{sec:debugging} عن تصحيح أخطاء الآلة من الخارج، ومن ذلك واجهات الدماغ والحاسوب. والفصلان~\ref{sec:eetypes} و\ref{sec:dendrites} يصنّفان العصبونات مرة أخرى، هذه المرة بحسب وظيفتها الكهربائية وعدد مداخلها. والفصل~\ref{sec:sleep} عمّا تفعله الآلة حين تنفصل عن العالم، والفصلان~\ref{sec:livingmachine} و\ref{sec:dishbrain} عن آلات مجمّعة من عصبونات حية على رقاقة.

\section{تمارين}

تمارين هذا الفصل تربط نتائج فصول مختلفة: كل تمرين يستند إلى فصلين على الأقل.

\begin{exercise}[\lvlA]
\label{ex:10:1}
\emph{الناقل والسجلات.} قنوات البث الأربع حزم من خمسة خطوط تقريبًا لكلٍّ منها (القضية~\ref{prop:bundle})؛ كل خط يتغير مرة في الثانية ويمكن تمييز أربعة مستويات عليه. كم سنة يحتاج التحكم لينقل ما ينقله ناقل \nt{GLUT} ($8.6\cdot10^{10}$ عصبون بالتردد~\eqref{eq:energy}، ودقة $1\ms$ بحسب~\eqref{eq:spikeentropy}) في ثانية واحدة؟
\end{exercise}

\begin{exercise}[\lvlB]
\label{ex:10:2}
\emph{مقبضان على حلقة واحدة.} في الحلقة~\eqref{eq:ei} تعمل مقابض~\eqref{eq:machine} على الإثارة: $\tau_E\dot E=-E+g[(1-a)w_{EE}E-w_{EI}I+h]$، ولا تتحكم في \nt{GABA}. بأرقام الشكل~\ref{fig:ei} ترفع رشقةُ \nt{NORA} الكسبَ $g$ إلى $6$. ما أصغر $a$ تكون عنده الحلقة مستقرة؟ وهل يمكن إدارة مقبضي \nt{NORA} و\nt{ACET} كلٍّ على حدة؟
\end{exercise}

\begin{exercise}[\lvlB]
\label{ex:10:3}
\emph{من أين يأتي ثمن البث.} مخارج $N$ عصبونًا $y_k=f(u_k)+\xi_k$ بضجيج غاوسي مستقل تباينه $\sigma^2$؛ و$R\approx R_0+\sum_kR'_k\xi_k$ بقيم $R'_k$ متساوية، ويبث \nt{DOPA} القيمة $\delta=R-R_0$. أوجد نسبة الإشارة إلى الضجيج للتصحيح $\delta\,\xi_jx_j$ في محاولة واحدة. كيف يزداد عدد المحاولات مع $N$؟
\end{exercise}

\begin{exercise}[\lvlB]
\label{ex:10:4}
\emph{ربط الصوت بالومضة.} يصل الصوت إلى الناقل بعد $5\ms$، والومضة بعد $30\ms$ مضافًا إليها كمون النبضة الأولى~\eqref{eq:latency} ($\tau=10\ms$، و$U$ من $1.2\theta$ إلى $4\theta$)؛ والخلفية على كل مدخل $5$ نبضات في الثانية. اختر تأخير خط الصوت وأصغر نافذة تزامن تصلح لكل التباينات. كم ربطًا كاذبًا في الثانية ستعطي الخلفية؟
\end{exercise}

\begin{exercise}[\lvlB]
\label{ex:10:5}
\emph{نمذجة: من يلاحظ التغيير.} يرى الناقد $y=s+\text{ضجيج}$ ($R=1$)؛ وباحتمال $1/200$ تقفز $s$ إلى قيمة جديدة تباينها $J^2=4$. نمذج مرشح كالمان بـ$q=0$ يرفع $P$ إلى $J^2$ حين تتجاوز $E\leftarrow0.9E+\delta^2/(P+R)$ القيمة $20$. بكم يقل خطؤه عن خطأ أفضل $\alpha$ ثابت؟
\end{exercise}

\begin{exercise}[\lvlA]
\label{ex:10:6}
\emph{كم محاولة لتغيير العادة.} تعلّمت الشبكة $Q_A=0.8$ و$Q_B=0.2$، وتختار بحسب~\eqref{eq:softmax} مع $\sigma=1$ و$g=8$. الآن يعطي $A$ العصير باحتمال $0.2$؛ $Q_A\leftarrow Q_A+\alpha(r-Q_A)$، و$Q_B$ لا يتغير. بعد كم اختيارًا لـ$A$ ينخفض احتمال اختيار $A$ إلى $0.6$ عند $\alpha=0.1$ و$0.5$؟ وماذا لو خفّض \nt{NORA} الكسب $g$ إلى $2$؟
\end{exercise}

\begin{exercise}[\lvlC]
\label{ex:10:7}
\emph{حزمة أسلاك بدل سلك.} المخارج $y_k=w_k+\xi_k$، والمكافأة $R=-\sum_ky_k^2$؛ يبث كل خط من خطوط الحزمة إلى مجموعة من $G$ عصبونًا خطأَ مخارجها، $w_k\leftarrow w_k+\eta\,\delta_l\xi_k$. بيّن أنه عند أفضل $\eta$ يلزم $\approx(G+2)\ln(1/\varepsilon)$ محاولة للوصول إلى $\sum w_k^2=\varepsilon\sum w_k^2(0)$. كم تتعلم دارة من $10^6$ عصبون على خمسة خطوط عند $\varepsilon=0.01$، بمحاولة كل ثلاث ثوانٍ؟
\end{exercise}
